\documentclass[10pt,a4paper]{amsart}
\usepackage[margin=19mm]{geometry}
\usepackage{amsmath,amssymb,mathtools}
\usepackage[scr=rsfs,cal=euler]{mathalfa}
\usepackage{xfrac}
\usepackage[unicode,hidelinks,bookmarks=false]{hyperref}
\newcommand{\ie}{i.e.\ }
\newcommand{\cf}{cf.\ }
\newcommand{\resp}{respectively\ }
\newcommand{\Subsection}{\subsection}
\providecommand{\nobreakdash}{\nobreak\hbox{-}}
\setlength{\emergencystretch}{2em}
\multlinegap0pt
\usepackage{amssymb}
\usepackage{siunitx}
\usepackage[shortlabels]{enumitem}
\usepackage{mathtools}
\usepackage[normalem]{ulem}
\usepackage{tcolorbox}
\newtheorem{thm}{Theorem}
\newtheorem{lem}{Lemma}
\newcommand{\Bea}{\begin{eqnarray*}}
	\newcommand{\Eea}{\end{eqnarray*}}
\title{Continuous family of compact quantum metric space structures from cocycle twisted crossed product $\textrm{C}^{\ast}$-algebras}
\author{Arnab Chattopadhyay, Soumalya Joardar}
\date{Source: arXiv:2609.11316v1 (CC BY 4.0)}
\begin{document}
\maketitle

\noindent\emph{Source and licence.} Continuous family of compact quantum metric space structures from cocycle twisted crossed product $\textrm{C}^{\ast}$-algebras. Version arXiv:2609.11316v1 (CC BY 4.0), distributed by arXiv under the Creative Commons Attribution 4.0 International licence, http://creativecommons.org/licenses/by/4.0/, as recorded in the arXiv licence metadata for this exact version and shown on the abstract page https://arxiv.org/abs/2609.11316v1. Full licence notice: \texttt{licenses/arxiv-2609.11316-CC-BY-4.0.txt}.\par\smallskip

\noindent\textit{Editorial scope.} Counted unit: the article's thm thm (source0.tex lines 718--720). The article's complete original proof of it is delivered with it (source0.tex lines 722--774, one \texttt{proof} environment of 53 lines). No mathematics is written, repaired or re-derived: the statement and the proof are the article's own lines, in the article's own notation, and no retained line is substituted or re-flowed.  Retained uncounted as the source-local context the proof needs: lines 424--426; lines 695--695; lines 1163--1176.  Nothing was omitted from any retained span, so the package carries no omission record.  No retained line contains a citation, so the wrapper carries no bibliography and no bibliography entry.  No retained line contains a figure environment, an image-inclusion command or drawing code, so no image is delivered and the package declares no asset dependency.  Editorial formatting only, in addition: an AMS \texttt{amsart} preamble that restates the article's own declarations and macros used by the retained lines, namely the environments \texttt{thm}, \texttt{lem} on separate counters, together with the standard packages those lines need.  The retained environments print in this document as Theorem 1, Lemma 1 in order of appearance, whereas the article prints the same items with the numbering of its own full document.  The complete native source of the article as distributed in the arXiv e-print for this exact version is bundled unchanged as \texttt{sources/210007/source0.tex}.
\begin{thm}\label{totbouiff}
    Let $L$ be a Lipschitz seminorm on an operator system $X$ and let $\sigma\in S(X)$. Then $L$ is a Lip-norm if and only if the set $\{x\in X: L(x)\leq 1, \sigma(x)=0\}$ is norm totally bounded in X. 
\end{thm}

\subsection{The family of compact quantum metric spaces} \label{Lip-norm defn}

\begin{lem} \label{threshold}
Let $\alpha, \beta, \gamma,$ and $C$ be real constants satisfying $\beta \ge \alpha > 0$ and $\gamma > C > 0$, and define $\eta = \gamma - C$. For any $\varepsilon > 0$, there exists an integer $N \ge 1$ such that for all integers $n \ge N$,
\begin{equation*}
    \sum_{k=n+1}^{\infty} e^{-\gamma k^\beta + C k^\alpha} \le \varepsilon.
\end{equation*}
Furthermore, an explicit such threshold $N$ is given by
\begin{equation*}
    N = \begin{cases}
    \left\lceil \left(\frac {1} {\eta} \log \left (\frac {1} {\eta \beta \varepsilon} \right) \right )^{\frac{1}{\beta}} \right \rceil, & \mathrm{if }\ \beta \ge 1, \\[12pt]
    \left\lceil \left( \frac{2}{\eta} \log \left (\frac {K} {\varepsilon} \right) \right)^{\frac{1}{\beta}} \right\rceil, & \mathrm{if }\ 0 < \beta < 1,
    \end{cases}
\end{equation*}
where $K = \frac{2}{\eta \beta} \left( \frac{2(1-\beta)}{e \eta \beta} \right)^{\frac{1-\beta}{\beta}}.$
\end{lem}

\begin{thm}
Let $(A,\rho,\Gamma)$ be a $C^{\ast}$-dynamical system satisfying the standard assumptions; $\Gamma$ admits a normalized $2$-cocycle $\sigma : \Gamma \times \Gamma \rightarrow \mathbb{T}$. If $\rho$ is $\mathrm{Lip}$-$\mathrm {isometric}$ then for all $\beta \geq \alpha$ and $\gamma > C,$ the associated reduced twisted crossed product $C^{\ast}$-algebra $A_{\sigma} : = A \rtimes_{r, \rho, \sigma} \Gamma$ has a CQMS structure with respect to the stratified Lip-norm $L^{S, \beta, \gamma}_{\ell}$ as defined in Subsection \ref{Lip-norm defn}.
\end{thm}

\begin{proof} Let $\varepsilon > 0.$ First we get hold of some threshold $N \in \mathbb N$ such that for all $n \geq N$ and for all $\mathfrak{a} = \sum\limits_{t \in \Gamma} \delta_t a_t \in C_c (\Gamma, \mathcal A)$ with $L_{\ell}^{S, \beta, \gamma} (\mathfrak{a}) \leq 1,$ we have 
$$\left \|\sum\limits_{\ell (t) \geq n + 1} \delta_t a_t \right \|_{\mathrm {red}} < \frac {\varepsilon} {2}.$$
First note that if $L_{\ell}^{S, \beta, \gamma} (\mathfrak{a}) \leq 1,$ then for all $t \in \Gamma \setminus \{e\},$ we have 
$$\left \|a_t \right \| \leq e^{- \gamma \ell (t)^{\beta}}.$$
Now fix some $n \geq 1.$ Then we have
\Bea
\left \|\sum\limits_{\ell (t) \geq n + 1} \delta_t a_t \right \|_{\mathrm {red}} & \leq & \sum\limits_{\ell (t) \geq n + 1} \left \|a_t \right \| \\ 
& \leq & \sum\limits_{\ell (t) \geq n + 1} e^{- \gamma \ell (t)^{\beta}} \\ 
& \leq & \sum\limits_{k = n + 1}^{\infty} \sum\limits_{k - 1 < \ell (t) \leq k} e^{- \gamma k^{\beta}} \\ 
& \leq & \sum\limits_{k = n + 1}^{\infty} e^{- \gamma k^{\beta}} e^{C k^{\alpha}}.
\Eea
By Lemma \ref{threshold}, it then follows that the right hand side of the above inequality is less than $\frac {\varepsilon} {2}$ for sufficiently large $n.$ Let $N$ be this threshold.

\vspace{2mm}

Now fix some state $\psi$ of $A,$ and define a state $\tau$ on the reduced crossed product $A_{\sigma} = A \rtimes_{r, \rho, \sigma} \Gamma$ by $\tau(\mathfrak{a}) = \psi(E(\mathfrak{a})),$ where $E : A_{\sigma} \rightarrow A$ is the canonical conditional expectation. Thus, for $\mathfrak{a} = \sum\limits_{t \in \Gamma} \delta_t a_t,$ we have $\tau(\mathfrak{a}) = \psi (a_e).$

To establish the CQMS structure on $A_{\sigma},$ we must show that the set 
$$\widetilde{\mathcal{L}}_1^{S, \beta, \gamma} := \left \{\mathfrak{a} \in A_{\sigma}\ :\ L_{\ell}^{S, \beta, \gamma} (\mathfrak{a}) \leq 1 \ \mathrm{and}\ \tau(\mathfrak{a}) = 0 \right \}$$ 
is totally bounded with respect to the reduced norm. Consequently, it is enough to show that the set of finite truncations 
$$S := \left \{\mathfrak{a}=\sum\limits_{\ell (t) \leq N} \delta_t a_t\ :\ \mathfrak{a} \in \widetilde{\mathcal{L}}_1^{S, \beta, \gamma} \right \}$$ 
forms a totally bounded set. Note that for $\mathfrak{a} \in \widetilde{\mathcal{L}}_1^{S, \beta, \gamma},$ we automatically have $\psi(a_e) = \tau(\mathfrak{a}) = 0.$
Since $(A, L_A)$ is a CQMS, the set
$$\mathcal L_1' : = \left \{a \in A\ :\ L_A (a) \leq 1\ \mathrm {and}\ \psi (a) = 0 \right \}$$ 
is totally bounded in $A$ by Theorem \ref{totbouiff}. So we can get hold of $c_1, \cdots, c_p \in \mathcal L_1'$ such that 
$$\mathcal L_1' \subseteq \bigcup\limits_{i = 1}^{p} B \left (c_i, \frac {\varepsilon} {4 \lambda_N} \right ),$$ 
where $\lambda_N : = \left \lvert \left \{t \in \Gamma\ :\ \ell (t) \leq N \right \} \right \rvert.$

Furthermore, for $t \neq e,$ the condition $L_{\ell}^{S, \beta, \gamma}(\mathfrak{a}) \leq 1$ implies $\|a_t\| \leq e^{-\gamma \ell(t)^\beta} \leq 1.$ Therefore, the scalars $\psi (a_t)$ lie in the closed unit disk $D := \{z \in \mathbb{C} \ : \ |z| \leq 1\}.$ Since $D$ is compact, there exists a finite set $Z = \{z_1, \dots, z_m\} \subseteq D$ such that
$$D \subseteq \bigcup_{j=1}^m B\left(z_j, \frac{\varepsilon}{4 \lambda_N}\right).$$
Now, let $\mathfrak{a} \in \widetilde{\mathcal{L}}_1^{S, \beta, \gamma}.$ We construct an approximation for each $a_t$ with $\ell(t) \leq N.$ 
For $t = e,$ we have $L_A(a_e) \leq 1$ and $\psi (a_e) = 0.$ Thus $a_e \in \mathcal{L}_1',$ and we can choose $b_e \in \{c_1, \dots, c_p\}$ such that
$$\left \|a_e - b_e \right \| < \frac{\varepsilon}{4 \lambda_N}.$$
For $g \neq e,$ we have $L_A(a_t) \leq 1.$ Because $L_A$ vanishes on scalars, we have $L_A(a_t - \psi(a_t)1_A) = L_A(a_t) \leq 1,$ and clearly $\psi (a_t - \psi(a_t)1_A) = 0.$ Thus $a_t - \psi(a_t)1_A \in \mathcal{L}_1',$ and we can choose $b_t \in \{c_1, \cdots, c_p\}$ such that
$$\left \|a_t - \psi (a_t) 1_A - b_t \right \| < \frac {\varepsilon} {4 \lambda_N}.$$
Simultaneously, since $\psi(a_t) \in D,$ we can choose $w_t \in Z$ such that 
$$\left |\psi (a_t) - w_t \right | < \frac{\varepsilon}{4 \lambda_N}.$$
Define $y_t = b_t + w_t 1_A$ for $t \neq e,$ and $y_e = b_e.$ Notice that each $y_t$ belongs to the finite set
$$F := \{ c_i + z_j 1_A \ : \ 1 \leq i \leq p, 1 \leq j \leq m \} \cup \{c_1, \dots, c_p\}.$$ 
We then have, for $t \neq e$
\Bea
\left \| a_t - y_t \right \| & = & \left \| (a_t - \psi (a_t)1_A - b_g) + (\psi (a_t) - w_t)1_A \right \| \\
& \leq & \left \| a_t - \psi (a_t)1_A - b_t \right \| + \left | \psi (a_t) - w_t \right | \\
& < & \frac{\varepsilon}{4 \lambda_N} + \frac{\varepsilon}{4 \lambda_N} \\ & = & \frac{\varepsilon}{2 \lambda_N}.
\Eea
For $t = e,$ we already established $\|a_t - y_t\| < \frac{\varepsilon}{4 \lambda_N} < \frac {\varepsilon} {2 \lambda_N}.$
So we have
\Bea
\left \|\sum\limits_{\ell (t) \leq N} \delta_t a_t - \sum\limits_{\ell (t) \leq N} \delta_t y_t \right \|_{\mathrm{red}} & \leq & \sum\limits_{\ell (t) \leq N} \left \|a_t - y_t \right \| \\
& < & \sum\limits_{\ell (t) \leq N} \frac{\varepsilon}{2 \lambda_N} \\ & = & \frac{\varepsilon}{2}.
\Eea
This shows that the set $S$ is totally bounded with respect to the reduced norm, as required.
\end{proof}

\end{document}
